\section{Residual-Cost Extension and Adoption Thresholds}
\label{app:partial-erosion}

This appendix records the product-market and private-adoption algebra underlying Section~\ref{sec:selective-erosion}. The extension changes one primitive relative to the canonical full-bypass technology: after a firm adopts a foreign formal standard, it becomes fully compatible with firms using that standard but may retain residual marginal adaptation cost
\[
d=\lambda c,
\qquad
0\leq\lambda\leq1.
\]
Native support under the formal partition continues to carry zero adaptation cost.

\subsection{Member market after outsider adoption}

After the outsider adopts the bloc standard, all three firms are compatible in a member market. The two member firms have zero marginal cost and the outsider has marginal cost $d$. Let $x$ denote each member's output and $y$ the outsider's output. The compatible-firm first-order conditions reduce to
\[
1-4(1-v)x+d=0
\quad\text{for the members,}
\]
together with the outsider condition implied by the common compatible-demand slope. Solving the linear system gives
\begin{equation}
x(d)=\frac{1+d}{4(1-v)},
\qquad
y(d)=\frac{1-3d}{4(1-v)}.
\label{eq:appendix-partial-quantities}
\end{equation}
Operating profits and consumer surplus are therefore
\begin{align}
A_R(d)&=\frac{(1+d)^2}{16(1-v)},\
B_R(d)&=\frac{(1-3d)^2}{16(1-v)},\
K_R(d)&=\frac{(3-d)^2}{32(1-v)^2}.
\label{eq:appendix-partial-blocks}
\end{align}
At $d=0$, these reduce to $A_R(0)=B_R(0)=P$ and $K_R(0)=K_I$.

Substitution into
\[
W_M^{SU,O}(d)=K_R(d)+2A_R(d)+C
\]
gives
\begin{equation}
W_M^{SU,O}(d)-W^{IS}
=
-\mathscr D+
\frac{d[(5-4v)d+2(1-4v)]}{32(1-v)^2},
\label{eq:appendix-residual-identity}
\end{equation}
which is equation~\eqref{eq:partial-erosion-identity}.

\subsection{Outsider adoption threshold}

Before adoption, the outsider earns $B$ in each member market. After adoption with residual cost $d$, it earns $B_R(d)$. Define
\begin{equation}
T_O(d)=B_R(d)-B.
\label{eq:appendix-TOd}
\end{equation}
Because one bloc-standard investment covers both member markets, the symmetric outsider strictly adopts iff
\begin{equation}
F<2T_O(d).
\label{eq:appendix-outsider-partial-threshold}
\end{equation}
At $d=0$, $T_O(0)=P-B=T_A$.

\subsection{Reverse adoption by one SU member}

Now consider the outsider market. Suppose one SU member adopts the outsider standard and bears residual marginal cost $d$, while the other member does not adopt and continues to bear cost $c$. Let $q_N(d)$ be the native outsider's quantity, $q_A(d)$ the adopting member's quantity, and $q_S(d)$ the non-adopting member's quantity. The unique interior Cournot solution is
\begin{align}
q_N(d)
&=
\frac{(1-v)(1+c)+d(1-2v)}
     {2(1-v)(2-3v)},\
q_A(d)
&=
\frac{(1-v)(1+c)-d(3-4v)}
     {2(1-v)(2-3v)},\
q_S(d)
&=
\frac{1-3v-3c(1-v)+d}
     {2(2-3v)}.
\label{eq:appendix-one-adopter-quantities}
\end{align}
The adopting member earns
\[
A_A(d)=(1-v)q_A(d)^2.
\]
Hence the unilateral reverse-adoption threshold under SU is
\begin{equation}
T_U(d)=A_A(d)-C.
\label{eq:appendix-TUd}
\end{equation}
Starting from separate national standards, the same post-adoption state is reached but the foreign firm's pre-adoption profit is $S$, so
\begin{equation}
T_W(d)=A_A(d)-S.
\label{eq:appendix-TWd}
\end{equation}
Consequently,
\begin{equation}
T_W(d)-T_U(d)=C-S>0
\label{eq:appendix-TW-TU-partial}
\end{equation}
on the canonical domain.

\subsection{Adoption after the other member has adopted}

If both SU members adopt the outsider standard, the two adopted members each bear residual cost $d$ and all three firms are compatible. Each adopted member produces
\begin{equation}
q_{AA}(d)=\frac{1-2d}{4(1-v)}
\end{equation}
and earns
\begin{equation}
A_{AA}(d)=\frac{(1-2d)^2}{16(1-v)}.
\end{equation}
Before the second member adopts, it is the singleton in the one-adopter state and earns $q_S(d)^2$. Its post-rival adoption threshold is therefore
\begin{equation}
T_A^R(d)=A_{AA}(d)-q_S(d)^2.
\label{eq:appendix-TARd}
\end{equation}
At $d=0$,
\[
T_U(0)=T_U,
\qquad
T_W(0)=T_W,
\qquad
T_A^R(0)=T_A.
\]

\subsection{Selection-free interval and exact witness}

Define
\[
F_L(d)=\max\{0,T_W(d),T_A^R(d)\},
\qquad
F^*(d)=2T_O(d).
\]
If $F_L(d)<F<F^*(d)$, outsider adoption is strictly optimal and both member non-adoption and SW non-adoption are strict against every relevant rival action.

At
\[
(c,v,\lambda,F)
=
\left(
\frac1{10},
\frac6{25},
\frac12,
\frac9{100}
\right),
\qquad
d=\frac1{20},
\]
the exact thresholds are
\begin{align}
T_U(d)&=\frac{42273}{1245184},\
T_W(d)&=\frac{2122041}{31129600},\
T_A^R(d)&=\frac{2024181}{31129600},\
F^*(d)&=\frac{459189}{3891200}.
\end{align}
Thus
\[
F_L(d)
=
\frac{2122041}{31129600}
<
\frac9{100}
<
\frac{459189}{3891200}
=
F^*(d).
\]
Together with the strict pre- and post-adaptation welfare inequalities reported in Proposition~\ref{prop:incomplete-adaptation-open-set}, this establishes the nonempty joint open set used in the main text.
